
\section{Threat Model and Security Properties}\label{sec:sys_model}

In this section, we present our threat model and define the privacy and the desired security properties of our system.
As illustrated in Fig.~\ref{fig:architecture}, the system involves three types of entities: clients, a set of signing authorities, and network nodes. 
Our construction relies on a threshold multi-authority model, where authorities hold shares of a distributed signing key, and a client combines partial credentials from a sufficient subset of authorities. 
Then, this credential is re-randomized and bound to the encrypted destination for each packet. 
Any mixnode along the route can verify the credential without learning the source or the destination.

Signing authorities and mixnodes are assumed to be \emph{honest-but-curious}: they correctly execute the protocol but are not trusted for privacy guarantees. 
Security is ensured as long as fewer than a threshold $t$ of signing authorities collude.

\subsection{Threat Model}

We consider cheating client, corrupted mixnodes and authorities, as well as global passive and active network adversaries. 
We assume the hardness of the Elliptic Curve Discrete Logarithm Problem (ECDLP) in the underlying group.
Hash functions are modeled as random oracles, and HMAC is assumed to be existentially unforgeable under chosen-message attacks.

\smallskip
\noindent{\textbf{Cheating Client.}}
A malicious client may attempt to reach a destination for which it does not have credential.

\smallskip
\noindent{\textbf{Corrupted Mixode.}}
A corrupted mixnode may reveal its entire internal state, including private keys and shared secrets.
Compromising a node must not compromise the privacy or integrity of the remaining path, and unlinkability is preserved as long as at least one node on the path is honest.

\smallskip
\noindent{\textbf{Corrupted Authorities.}}
Corrupted authorities may attempt to infer additional information from their computation.
We allow collusion among authorities, provided that fewer than $t$ authorities are colluding together.
Under this assumption, colluding authorities cannot reconstruct the distirbuted signing key.

\smallskip
\noindent{\textbf{Passive Network Adversary.}}
A passive adversary can observe traffic at any point in the network but cannot modify, inject, or drop packets.
Its goal is to break unlinkability by correlating incoming and outgoing packets.

\smallskip
\noindent{\textbf{Active Network Adversary.}}
An active adversary can arbitrarily modify, drop, replay, or inject packets.
Its objective is to forge valid integrity tags to tamper with headers while remaining undetected. 



\subsection{Security and Privacy Objectives}\label{sec:sp-objectives}


Our DSphinx construction aims to preserve the security guarantees of Sphinx~\cite{sphinx}, including unlinkability and integrity protection, while extending them to a decentralized route sampling setting.

\smallskip
\noindent\textbf{Unlinkability.} We consider \emph{node-level unlinkability}, requiring that an adversary cannot link an incoming packet to its corresponding outgoing packet at any node. This is achieved if all per-hop header components are computationally unlinkable without knowledge of the node's secret key.

\smallskip
\noindent{\textbf{Integrity-Forgery Resistance.}}
As in Sphinx, DSphinx employs per-hop integrity tags to ensure that any modification of packet headers can be detected by honest nodes.
An active adversary must not be able to alter a packet in transit while producing a valid integrity tag.

% \smallskip
% \noindent{\textbf{Routing Sampling Security.}}
% The route sampling procedure must remain secure against malicious clients and corrupted STTPs.
% In particular, no adversary or colluding set of adversaries should be able to bias, predict, or influence the route selection process beyond the intended randomness of the protocol.

\smallskip
\noindent{\textbf{Collusion Resistance.}} 
We require that collusions between clients, STTPs, and network nodes reveal no additional routing or anonymity information beyond their individual views. 
In particular, as long as fewer than $d$ STTPs are corrupted, adversaries cannot reconstruct complete routing information or compromise the security of honest participants.
